% =====================================================================
%  Fundación Llaftun — Lo que necesito de ustedes
%  Documento de pendientes del sitio web
%
%  Compilar con pdfLaTeX:   pdflatex pendientes-llaftun.tex
%  (dos veces, para que hyperref resuelva los enlaces)
%
%  Tipografía: Alegreya y Alegreya Sans, las mismas del sitio.
%  Si tu distribución no las trae, comenta las dos líneas marcadas
%  con [FUENTES] y el documento compila igual con la tipografía por
%  defecto.
% =====================================================================

\documentclass[11pt,a4paper]{article}

\usepackage[T1]{fontenc}
\usepackage[utf8]{inputenc}
\usepackage[spanish,es-noquoting,es-tabla]{babel}

% [FUENTES] — comentar estas dos líneas si no están instaladas
\usepackage{Alegreya}
\usepackage{AlegreyaSans}

\usepackage[
  a4paper,
  top=2.4cm, bottom=2.4cm, left=2.6cm, right=2.6cm
]{geometry}
\usepackage{xcolor}
\usepackage{enumitem}
\usepackage{tcolorbox}
\usepackage{titlesec}
\usepackage{microtype}
\usepackage[hidelinks]{hyperref}

% ---------------------------------------------------------------------
%  Colores: los mismos del sitio, tomados del logotipo
% ---------------------------------------------------------------------
\definecolor{bosque}{HTML}{2D5843}      % verde de los títulos
\definecolor{bosqueclaro}{HTML}{4D7A61} % verde del antetítulo
\definecolor{greda}{HTML}{9C4230}       % rojo tierra: urgencias y datos
\definecolor{gredasuave}{HTML}{F6ECE9}
\definecolor{pizarra}{HTML}{3F5A6B}     % azul: decisiones
\definecolor{pizarrasuave}{HTML}{EAF0F3}
\definecolor{tinta}{HTML}{1B241E}
\definecolor{tintasuave}{HTML}{4B5545}
\definecolor{superficie}{HTML}{F0F3EE}
\definecolor{superficiealt}{HTML}{E6EBE3}

\color{tinta}

% ---------------------------------------------------------------------
%  Títulos de sección
% ---------------------------------------------------------------------
\titleformat{\section}
  {\normalfont\Large\bfseries\color{bosque}}
  {}{0pt}{}
  [\vspace{2pt}{\color{superficiealt}\titlerule[1pt]}]
\titlespacing*{\section}{0pt}{2.2em}{0.8em}

\setlength{\parindent}{0pt}
\setlength{\parskip}{0.5em}

% ---------------------------------------------------------------------
%  Etiquetas de tipo (los "chips" del documento web)
% ---------------------------------------------------------------------
\newcommand{\etiqueta}[3]{%
  \,\raisebox{0.12ex}{\colorbox{#1}{%
    \footnotesize\sffamily\color{#2}\,#3\,}}%
}
\newcommand{\contenido}{\etiqueta{superficiealt}{tintasuave}{CONTENIDO}}
\newcommand{\decision}{\etiqueta{pizarrasuave}{pizarra}{DECISIÓN}}
\newcommand{\dato}{\etiqueta{gredasuave}{greda}{DATO}}

% ---------------------------------------------------------------------
%  Lista de puntos: código a la izquierda, título y cuerpo a la derecha
% ---------------------------------------------------------------------
\newlist{puntos}{description}{1}
\setlist[puntos]{
  style=sameline,
  leftmargin=4.6em,
  labelwidth=3.9em,
  labelsep=0.7em,
  align=left,
  font=\normalfont\sffamily\small\color{tintasuave},
  itemsep=1.2em,
  parsep=0.45em,
  topsep=1.1em
}

% Título de un punto
\newcommand{\punto}[1]{\textbf{#1}}

% Antetítulo de sección (la bajada explicativa bajo el título)
\newcommand{\nota}[1]{{\small\color{tintasuave}#1}\par\vspace{0.2em}}

% ---------------------------------------------------------------------
\begin{document}
\sffamily

% =====================================================================
%  Portada
% =====================================================================
{\footnotesize\color{bosqueclaro}FUNDACIÓN LLAFTUN \textperiodcentered{} SITIO WEB NUEVO}

\vspace{0.7em}

{\rmfamily\bfseries\color{tinta}\fontsize{30}{34}\selectfont Lo que necesito de ustedes\par}

\vspace{0.7em}

{\rmfamily\large\color{tintasuave}
Ya construí la estructura del sitio y rescaté los textos que estaban
publicados en la web actual. Lo que sigue es completar el material que
todavía no existe en el Drive.\par}

\vspace{1.2em}

\begin{tcolorbox}[
  colback=superficie, colframe=superficie,
  boxrule=0pt, arc=2pt, left=10pt, right=10pt, top=8pt, bottom=8pt
]
\small\color{tintasuave}
\textbf{\color{tinta}11} cosas por entregar \hfill
\textbf{\color{tinta}4} decisiones por tomar \hfill
\textbf{\color{tinta}5} datos técnicos \hfill
\textbf{\color{tinta}1} asunto urgente
\end{tcolorbox}

% =====================================================================
%  Urgente
% =====================================================================
\vspace{0.8em}

\begin{tcolorbox}[
  colback=gredasuave, colframe=greda,
  boxrule=0pt, leftrule=2.5pt, arc=0pt,
  left=12pt, right=12pt, top=10pt, bottom=10pt
]
{\footnotesize\bfseries\color{greda}ATENDER HOY}

\vspace{0.4em}
\small
En la carpeta \textbf{ACCESOS y LOGOS} del Drive hay un documento con la
contraseña del administrador de WordPress escrita en texto plano. Cualquier
persona con acceso a ese Drive puede entrar a modificar el sitio.

\vspace{0.5em}
Hay que cambiar esa contraseña por una nueva, borrar el documento, y pasar
los accesos por otro medio. No es urgente por el sitio nuevo: es urgente por
el que está publicado ahora.
\end{tcolorbox}

% =====================================================================
\section{Página de inicio}
\nota{Es la primera pantalla que ve quien llega al sitio: una frase, un
carrusel con los proyectos y un enlace a cada uno.}

\begin{puntos}

\item[IN\textperiodcentered1] \punto{La frase representativa}\contenido

Una sola frase, de una línea, que abra el sitio y diga quiénes son. Es lo
primero y lo más grande que se lee. Mientras no la tenga uso ``Entregamos
herramientas para el bienestar de quienes más lo necesitan'', tomada de la
misión, pero conviene que la escriban ustedes.

\item[IN\textperiodcentered2] \punto{Cuáles son los proyectos más representativos}\decision

Hay más de veinte proyectos documentados entre la web y el Drive. Elijan
entre cuatro y seis para el carrusel y la portada. Por ahora están puestos:
Contacto con la biodiversidad de mi territorio, El Tesoro del Bosque
Esclerófilo, Cooperativa de Trabajo Inclusiva e Hipoterapia.

\item[IN\textperiodcentered3] \punto{Una fotografía por cada proyecto elegido}\contenido

Horizontal y de buena calidad, porque se muestra a pantalla completa en el
carrusel. Ver las condiciones al final de esta página.

\item[IN\textperiodcentered4] \punto{Texto de ``El Tesoro del Bosque Esclerófilo''}\contenido

De este proyecto solo tengo el título. Falta el texto que describa en qué
consistió, como el que ya existe para el proyecto FPA 2025. Dos o tres
párrafos bastan.

\end{puntos}

% =====================================================================
\section{Quiénes somos}
\nota{Los principios, la misión y la visión ya están listos: se rescataron
del sitio actual. Falta el equipo.}

\begin{puntos}

\item[QS\textperiodcentered1] \punto{Descripción de David Uribe y Rodrigo Palacios}\contenido

Sus carpetas en el Drive están vacías. Tres o cuatro líneas de cada uno, en
el mismo tono que las de Claudia y Pilar, que ya están escritas y quedaron
muy bien.

\item[QS\textperiodcentered2] \punto{Retratos de los integrantes del directorio}\contenido

Sirve cualquier foto donde se les vea de frente y con luz pareja; no hace
falta un estudio. Lo importante es que las cinco se vean parecidas entre sí,
y que el recorte sea cuadrado con la cara centrada, porque se muestran en
círculo.

\item[QS\textperiodcentered3] \punto{Una imagen de apertura}\contenido

Horizontal, de una charla o de una actividad con gente, para abrir la página
antes del texto de principios.

\end{puntos}

% =====================================================================
\section{Áreas de acción}
\nota{Esta es la parte con el vacío más grande, y donde primero necesito que
me respondan.}

\begin{puntos}

\item[AA\textperiodcentered1] \punto{¿Educación Ambiental y Bienestar reemplazan a los cuatro caminos?}\decision

El sitio actual organiza todo el trabajo de la fundación en cuatro caminos:
Medio Ambiente, Artes, Salud Natural y Economía Social. El documento de
requerimientos habla de dos áreas: Educación Ambiental y Bienestar.

No es un cambio de diseño, es un cambio de identidad: buena parte de los
proyectos históricos son de arteterapia y de economía social, y hoy no
tendrían dónde ubicarse. ¿Es un reenfoque deliberado, o las dos áreas nuevas
conviven con lo anterior?

\item[AA\textperiodcentered2] \punto{Todo el contenido del área Bienestar}\contenido

En el Drive no hay nada de esta área: ni texto ni fotografías. Todo el
material disponible es de educación ambiental. Necesito uno o dos párrafos
que expliquen en qué consiste el área, y una fotografía que la represente.

\item[AA\textperiodcentered3] \punto{Una fotografía para Educación Ambiental}\contenido

Las fotos del proyecto FPA 2024 que están en el Drive sirven bien para esto:
son las únicas tomadas con cámara.

\end{puntos}

% =====================================================================
\section{Material educativo}
\nota{Los siete documentos pedagógicos del Drive ya están incorporados y se
van a poder descargar directamente desde el sitio.}

\begin{puntos}

\item[ME\textperiodcentered1] \punto{Revisar la descripción de cada documento}\contenido

Escribí una descripción breve para cada uno de los siete materiales a partir
de su título, pero son provisorias. Les pido revisarlas y corregirlas: son
el texto que decide si alguien descarga el material o no.

\item[ME\textperiodcentered2] \punto{Una fotografía de estudiantes con los materiales}\contenido

Es la imagen que abre esta página, según lo pedido en el documento de
requerimientos.

\item[ME\textperiodcentered3] \punto{¿Existen los videos de la fundación?}\decision

El documento de requerimientos pide una sección de videos, pero no hay
ningún archivo de video en el Drive ni conozco un canal de YouTube de la
fundación.

Si existen en otro lado, los subo a YouTube y los enlazo. Si no existen
todavía, saco la sección y la agrego cuando los tengan: es preferible a
dejar un espacio vacío.

\end{puntos}

% =====================================================================
\section{Secciones que se mantienen}
\nota{El documento de requerimientos describe cuatro páginas, pero el sitio
publicado hoy tiene más. Ya construí también Servicios, Alianzas,
Transparencia y Colabora, con la misma estructura que tienen ahora, para que
no se pierda nada al cambiar de sitio. Solo necesito dos cosas de ustedes.}

\begin{puntos}

\item[SM\textperiodcentered1] \punto{Confirmar los datos bancarios}\dato

Copié el número de cuenta, el RUT y el correo de donaciones desde la página
actual, pero no tengo cómo verificar que sigan vigentes. Como por ahí entran
los aportes, un dígito equivocado significa donaciones perdidas sin que
nadie se entere.

Necesito que alguien los revise contra la cartola antes de publicar.

\item[SM\textperiodcentered2] \punto{Documentos de transparencia que falten}\contenido

La página queda armada con lo que hay entre el sitio actual y el Drive:
balances 2020 a 2023, memorias 2022 a 2024, estados de resultados y los
convenios del Fondo Vivir Mejor y del Fondo de Protección Ambiental.

Si hay documentos más recientes, o alguno que deba estar publicado y no esté
en esa lista, háganmelo saber. Vale la pena revisar qué obliga a publicar
cada convenio.

\end{puntos}

% =====================================================================
\section{Identidad visual}
\nota{En el Drive dejaron cuatro combinaciones de color. Hay que elegir una.}

\begin{puntos}

\item[IV\textperiodcentered1] \punto{Elegir la paleta de colores}\decision

Preparé una comparación donde las cuatro paletas están aplicadas al sitio
real, con la misma portada y las mismas tarjetas de proyecto, para que se vea
cómo queda cada una y no solo los cuadrados de color:
\href{https://claude.ai/code/artifact/e10ee75b-abc8-4572-9530-843c7e9d559a}%
{\color{bosque}\underline{ver las cuatro opciones}}.

Basta con responder con el número. El color se puede cambiar más adelante
sin rehacer el sitio, así que esta decisión no detiene el trabajo.

\end{puntos}

% =====================================================================
\section{Datos técnicos}
\nota{Esto no es contenido, pero sin ello no puedo publicar.}

\begin{puntos}

\item[DT\textperiodcentered1] \punto{Datos del hosting}\dato

Qué empresa lo provee, qué plan tienen contratado, y el acceso al panel de
administración. Necesito verificar que el servidor permita configurar las
direcciones del sitio nuevo, y cuánto espacio queda disponible.

\item[DT\textperiodcentered2] \punto{Tipo de cuenta de Instagram}\dato

Para mostrar las publicaciones automáticamente en el sitio, la cuenta debe
estar configurada como Empresa o Creador de contenido, no como cuenta
personal. Se revisa en la configuración de la aplicación.

\item[DT\textperiodcentered3] \punto{Los dos formularios de inscripción están rotos}\dato

Los botones ``Inscríbete aquí'' de socio donante y de voluntariado llevan a
la dirección de edición del formulario, no a la de respuesta. Quien entra
desde el sitio ve un aviso de permiso denegado y no puede inscribirse. Lleva
tiempo así.

Necesito los enlaces correctos, los que aparecen al apretar ``Enviar''
dentro de cada formulario de Google.

\item[DT\textperiodcentered4] \punto{Autorizaciones de imagen}\dato

Varias de las fotografías muestran a niñas, niños y jóvenes de las escuelas
con las que trabajan. Antes de publicarlas necesito confirmar que existen
las autorizaciones firmadas por sus apoderados.

\end{puntos}

% =====================================================================
%  Cómo enviar las fotografías
% =====================================================================
\vspace{1.6em}

\begin{tcolorbox}[
  colback=superficie, colframe=superficie,
  boxrule=0pt, arc=2pt, left=12pt, right=12pt, top=10pt, bottom=10pt
]
{\footnotesize\bfseries\color{tintasuave}CÓMO ENVIAR LAS FOTOGRAFÍAS}

\vspace{0.5em}
\small\color{tintasuave}
\begin{itemize}[leftmargin=1.1em, itemsep=0.5em, topsep=0pt]
  \item \textbf{\color{tinta}El archivo original, no una foto reenviada por
        WhatsApp.} WhatsApp reduce las imágenes a una fracción de su tamaño
        y esa pérdida no se recupera. La mayor parte de las fotos del Drive
        llegaron por ahí y se ven bien en el teléfono, pero se ven borrosas
        al ocupar el ancho de una pantalla de computador.
  \item \textbf{\color{tinta}Horizontales} para portadas y proyectos. Las
        verticales obligan a recortar y casi siempre se pierde lo importante.
  \item \textbf{\color{tinta}Al menos 1600 píxeles de ancho.} Las fotos
        tomadas con cámara, como las del proyecto FPA 2024, cumplen de sobra.
  \item \textbf{\color{tinta}Con el nombre de lo que muestran.} Basta con que
        digan algo como ``huerto escuela El Roble'' en vez de ``WhatsApp
        Image 2025-10-20''.
\end{itemize}
\end{tcolorbox}

% =====================================================================
%  Cierre
% =====================================================================
\vspace{1.6em}
{\color{superficiealt}\hrule height 1pt}
\vspace{1em}

{\small\color{tintasuave}
No hace falta que llegue todo junto. Cada cosa que entreguen reemplaza un
espacio marcado en el sitio, y se puede ir revisando por partes.

Lo que más conviene resolver primero son las decisiones, porque de ellas
depende cómo se ordena el resto del sitio.\par}

\end{document}
